\section{Homework 6}

\begin{exercise}
    Find Weyl groups of \(A_2\), \(B_2\), and \(G_2\).
\end{exercise}

\begin{proof}
    Let \(\Phi\) be an irreducible root system of rank \(2\), and let
    \[
        \Delta=\{\alpha,\beta\}
    \]
    be a base. The Weyl group \(W\) is generated by the two simple reflections
    \[
        s_\alpha,\qquad s_\beta.
    \]
    If the angle between \(\alpha\) and \(\beta\) is \(\theta\), then the angle between
    the reflecting hyperplanes \(P_\alpha\) and \(P_\beta\) is
    \[
        \pi-\theta.
    \]
    The product of two reflections in lines meeting at angle \(\varphi\) is a
    rotation through angle \(2\varphi\). Hence
    \[
        s_\alpha s_\beta
    \]
    is a rotation through angle
    \[
        2(\pi-\theta).
    \]
    Thus \(W\) is a dihedral group:
    \[
        W\cong \langle s_\alpha,s_\beta
        \mid s_\alpha^2=s_\beta^2=1,\ (s_\alpha s_\beta)^m=1\rangle,
    \]
    where
    \[
        m=\frac{\pi}{\pi-\theta}.
    \]
    
    For \(A_2\), the two simple roots have equal length and satisfy
    \[
        \frac{2(\alpha,\beta)}{(\alpha,\alpha)}=-1.
    \]
    Thus
    \[
        \cos\theta=\frac{(\alpha,\beta)}
        {\|\alpha\|\|\beta\|}=-\frac12,
    \]
    so
    \[
        \theta=\frac{2\pi}{3}.
    \]
    Therefore
    \[
        \pi-\theta=\frac{\pi}{3},
        \qquad
        m=3.
    \]
    Hence
    \[
        W(A_2)\cong D_3\cong S_3,
    \]
    where \(D_3\) denotes the dihedral group of order \(6\).
    
    For \(B_2\), choose simple roots
    \[
        \alpha=e_1-e_2,\qquad \beta=e_2.
    \]
    Then
    \[
        (\alpha,\alpha)=2,\qquad (\beta,\beta)=1,\qquad
        (\alpha,\beta)=-1.
    \]
    Hence
    \[
        \cos\theta
        =
        \frac{(\alpha,\beta)}{\|\alpha\|\|\beta\|}
        =
        -\frac{1}{\sqrt2},
    \]
    so
    \[
        \theta=\frac{3\pi}{4}.
    \]
    Thus
    \[
        \pi-\theta=\frac{\pi}{4},
        \qquad
        m=4.
    \]
    Therefore
    \[
        W(B_2)\cong D_4,
    \]
    the dihedral group of order \(8\). Equivalently,
    \[
        W(B_2)\cong (\mathbb Z/2\mathbb Z)^2\rtimes S_2,
    \]
    the group of signed permutations of \(\{e_1,e_2\}\).
    
    For \(G_2\), the Cartan integers satisfy
    \[
        \langle \alpha,\beta\rangle\langle \beta,\alpha\rangle=3.
    \]
    Equivalently,
    \[
        4\cos^2\theta=3.
    \]
    Since the angle between two simple roots is obtuse, we get
    \[
        \cos\theta=-\frac{\sqrt3}{2},
        \qquad
        \theta=\frac{5\pi}{6}.
    \]
    Hence
    \[
        \pi-\theta=\frac{\pi}{6},
        \qquad
        m=6.
    \]
    Therefore
    \[
        W(G_2)\cong D_6,
    \]
    the dihedral group of order \(12\).
    
    In summary,
    \[
        W(A_2)\cong D_3\cong S_3,\qquad
        W(B_2)\cong D_4,\qquad
        W(G_2)\cong D_6,
    \]
    where \(D_m\) denotes the dihedral group of order \(2m\).
    \end{proof}

Let \(\Delta\subset\mathbb{R}^n\) be a root system, and let \(B=\{\beta_1,\ldots,\beta_n\}\) be a base of \(\Delta\). The Cartan matrix is defined as
\[
    \left(
        \frac{2(\beta_i\mid\beta_j)}{(\beta_j\mid\beta_j)}
    \right)_{1\leq i,j\leq n}.
\]

\begin{exercise}
    Prove that the Cartan matrix of a root system is positive definite.
\end{exercise}

\begin{proof}
    Let
    \[
        \Delta=\{\beta_1,\ldots,\beta_n\}
    \]
    be a base of the root system \(\Phi\), and let
    \[
        C=(C_{ij})_{1\leq i,j\leq n},
        \qquad
        C_{ij}=\frac{2(\beta_i,\beta_j)}{(\beta_j,\beta_j)}.
    \]
    
    Since the Cartan matrix need not be symmetric, the precise statement is that
    \(C\) is symmetrizable positive definite. Define
    \[
        D=\operatorname{diag}
        \left(
            \frac{(\beta_1,\beta_1)}{2},
            \ldots,
            \frac{(\beta_n,\beta_n)}{2}
        \right).
    \]
    Then
    \[
        (CD)_{ij}
        =
        C_{ij}\frac{(\beta_j,\beta_j)}{2}
        =
        \frac{2(\beta_i,\beta_j)}{(\beta_j,\beta_j)}
        \frac{(\beta_j,\beta_j)}{2}
        =
        (\beta_i,\beta_j).
    \]
    Therefore
    \[
        CD=
        \bigl((\beta_i,\beta_j)\bigr)_{1\leq i,j\leq n},
    \]
    which is the Gram matrix of the basis \(\Delta\).
    
    Since \(\Delta\) is a basis of the Euclidean space \(E\), the Gram matrix
    \[
        G=\bigl((\beta_i,\beta_j)\bigr)
    \]
    is symmetric positive definite. Indeed, for any nonzero vector
    \[
        x=(x_1,\ldots,x_n)^t\in\mathbb R^n,
    \]
    we have
    \[
        x^tGx
        =
        \left(
            \sum_{i=1}^n x_i\beta_i,
            \sum_{j=1}^n x_j\beta_j
        \right)
        =
        \left\|
            \sum_{i=1}^n x_i\beta_i
        \right\|^2
        >0,
    \]
    because \(\beta_1,\ldots,\beta_n\) are linearly independent.
    
    Thus
    \[
        CD
    \]
    is symmetric positive definite, and \(D\) is a positive diagonal matrix.
    Equivalently,
    \[
        D^{-1/2}CD^{1/2}
        =
        D^{-1/2}(CD)D^{-1/2}
    \]
    is symmetric positive definite. Hence \(C\) is symmetrizable positive definite.
    
    This is the standard sense in which the Cartan matrix of a finite root system
    is positive definite.
    \end{proof}

\begin{exercise}
    Find all Cartan matrices of rank \(3\).
\end{exercise}

\begin{proof}
    We classify Cartan matrices of rank \(3\) up to simultaneous permutation of
    rows and columns, i.e. up to reordering of the base
    \[
        \Delta=\{\beta_1,\beta_2,\beta_3\}.
    \]
    With the convention
    \[
        C_{ij}=\frac{2(\beta_i,\beta_j)}{(\beta_j,\beta_j)},
    \]
    we have
    \[
        C_{ii}=2,
    \]
    and, for \(i\neq j\),
    \[
        C_{ij}\in\{0,-1,-2,-3\}.
    \]
    Moreover,
    \[
        C_{ij}=0
        \quad\Longleftrightarrow\quad
        C_{ji}=0.
    \]
    Set
    \[
        n_{ij}=C_{ij}C_{ji}.
    \]
    Then
    \[
        n_{ij}\in\{0,1,2,3\}.
    \]
    
    We attach to \(C\) its Dynkin graph: the vertices are \(1,2,3\), and vertices
    \(i,j\) are joined by \(n_{ij}\) edges. If \(n_{ij}=2\) or \(3\), the direction of
    the arrow records which of \(|C_{ij}|\) and \(|C_{ji}|\) is larger.
    
    The associated symmetric matrix has diagonal entries \(2\) and off-diagonal
    entries
    \[
        -\sqrt{n_{ij}}.
    \]
    It is positive definite. In rank \(3\), this condition strongly restricts the
    possible graphs.
    
    First, a triangle cannot occur. Indeed, if the three edge multiplicities are
    \(a,b,c\geq 1\), then the determinant of
    \[
        \begin{pmatrix}
            2&-\sqrt a&-\sqrt c\\
            -\sqrt a&2&-\sqrt b\\
            -\sqrt c&-\sqrt b&2
        \end{pmatrix}
    \]
    is
    \[
        8-2(a+b+c)-2\sqrt{abc}\leq 8-6-2=0,
    \]
    contradicting positive definiteness.
    
    Thus a connected rank \(3\) Dynkin graph is a path. If its two edge
    multiplicities are \(a,b\in\{1,2,3\}\), then the determinant is
    \[
        8-2a-2b=2(4-a-b).
    \]
    Positive definiteness gives
    \[
        a+b<4.
    \]
    Hence the only connected possibilities are
    \[
        (a,b)=(1,1)
        \qquad\text{or}\qquad
        (a,b)=(1,2).
    \]
    These give the irreducible types
    \[
        A_3,\qquad B_3,\qquad C_3.
    \]
    
    The disconnected possibilities are obtained by taking direct sums of
    irreducible root systems of total rank \(3\). Thus the possible decomposable
    types are
    \[
        A_1\oplus A_1\oplus A_1,\qquad
        A_1\oplus A_2,\qquad
        A_1\oplus B_2,\qquad
        A_1\oplus G_2.
    \]
    
    Therefore, up to simultaneous permutation of rows and columns, the rank
    \(3\) Cartan matrices are precisely the following seven matrices:
    \[
    A_1\oplus A_1\oplus A_1:
    \qquad
    \begin{pmatrix}
    2&0&0\\
    0&2&0\\
    0&0&2
    \end{pmatrix}.
    \]
    \[
    A_1\oplus A_2:
    \qquad
    \begin{pmatrix}
    2&0&0\\
    0&2&-1\\
    0&-1&2
    \end{pmatrix}.
    \]
    \[
    A_1\oplus B_2:
    \qquad
    \begin{pmatrix}
    2&0&0\\
    0&2&-2\\
    0&-1&2
    \end{pmatrix}.
    \]
    \[
    A_1\oplus G_2:
    \qquad
    \begin{pmatrix}
    2&0&0\\
    0&2&-1\\
    0&-3&2
    \end{pmatrix}.
    \]
    \[
    A_3:
    \qquad
    \begin{pmatrix}
    2&-1&0\\
    -1&2&-1\\
    0&-1&2
    \end{pmatrix}.
    \]
    \[
    B_3:
    \qquad
    \begin{pmatrix}
    2&-1&0\\
    -1&2&-2\\
    0&-1&2
    \end{pmatrix}.
    \]
    \[
    C_3:
    \qquad
    \begin{pmatrix}
    2&-1&0\\
    -1&2&-1\\
    0&-2&2
    \end{pmatrix}.
    \]
    
    All other labelled rank \(3\) Cartan matrices are obtained from these by
    simultaneously permuting rows and columns, corresponding to a reordering of
    the simple roots.
    \end{proof}

\begin{exercise}
    Show that a Cartan matrix determines the root system uniquely up to isomorphism.
\end{exercise}

\begin{proof}
    Let \(\Phi\subset E\) and \(\Phi'\subset E'\) be two root systems with bases
    \[
        \Delta=\{\alpha_1,\ldots,\alpha_\ell\},
        \qquad
        \Delta'=\{\alpha_1',\ldots,\alpha_\ell'\}.
    \]
    Assume that the two Cartan matrices are the same, i.e.
    \[
        \frac{2(\alpha_i,\alpha_j)}{(\alpha_j,\alpha_j)}
        =
        \frac{2(\alpha_i',\alpha_j')}{(\alpha_j',\alpha_j')}
        \qquad
        \text{for all }1\leq i,j\leq \ell.
    \]
    Write
    \[
        C_{ij}
        =
        \frac{2(\alpha_i,\alpha_j)}{(\alpha_j,\alpha_j)}
        =
        \frac{2(\alpha_i',\alpha_j')}{(\alpha_j',\alpha_j')}.
    \]
    
    Since \(\Delta\) and \(\Delta'\) are bases of \(E\) and \(E'\), there is a unique
    linear isomorphism
    \[
        \varphi:E\longrightarrow E'
    \]
    such that
    \[
        \varphi(\alpha_i)=\alpha_i'
        \qquad
        \text{for }1\leq i\leq \ell.
    \]
    
    Let \(s_j=s_{\alpha_j}\) and \(s_j'=s_{\alpha_j'}\) be the corresponding simple
    reflections. For a simple root \(\alpha_i\), we have
    \[
        s_j(\alpha_i)
        =
        \alpha_i
        -
        \frac{2(\alpha_i,\alpha_j)}{(\alpha_j,\alpha_j)}
        \alpha_j
        =
        \alpha_i-C_{ij}\alpha_j.
    \]
    Similarly,
    \[
        s_j'(\alpha_i')
        =
        \alpha_i'-C_{ij}\alpha_j'.
    \]
    Therefore
    \[
    \begin{aligned}
        \varphi(s_j(\alpha_i))
        &=
        \varphi(\alpha_i-C_{ij}\alpha_j)        \\
        &=
        \alpha_i'-C_{ij}\alpha_j'               \\
        &=
        s_j'(\alpha_i')                         \\
        &=
        s_j'(\varphi(\alpha_i)).
    \end{aligned}
    \]
    Since the \(\alpha_i\) form a basis of \(E\), it follows that
    \[
        \varphi s_j=s_j'\varphi
        \qquad
        \text{for every }j.
    \]
    Hence \(\varphi\) conjugates the Weyl group \(W\) of \(\Phi\) onto the Weyl
    group \(W'\) of \(\Phi'\):
    \[
        \varphi W\varphi^{-1}=W'.
    \]
    
    We now use the standard fact that every root is Weyl-conjugate to a simple
    root. Indeed, if \(\gamma\in\Phi^+\) is not simple, then there exists some
    simple root \(\alpha_j\) such that
    \[
        (\gamma,\alpha_j)>0.
    \]
    Then \(s_j(\gamma)\) is again a positive root but has smaller height. Repeating
    this process sends \(\gamma\) to a simple root. Negative roots are handled by
    multiplying by \(-1\).
    
    Thus every root of \(\Phi\) has the form
    \[
        w(\alpha_i)
    \]
    for some \(w\in W\) and some \(i\). Therefore
    \[
    \begin{aligned}
        \varphi(w(\alpha_i))
        &=
        (\varphi w\varphi^{-1})(\varphi(\alpha_i)) \\
        &=
        w'(\alpha_i')
    \end{aligned}
    \]
    for some \(w'\in W'\). Hence
    \[
        \varphi(\Phi)\subseteq \Phi'.
    \]
    Applying the same argument to \(\varphi^{-1}\), we get the reverse inclusion,
    so
    \[
        \varphi(\Phi)=\Phi'.
    \]
    
    Moreover, the Cartan integers are preserved. Indeed, for roots obtained from
    simple roots by Weyl group action, preservation follows from the conjugation
    relation
    \[
        \varphi s_\alpha\varphi^{-1}=s_{\varphi(\alpha)}.
    \]
    Thus
    \[
        \left\langle \varphi(\beta),\varphi(\alpha)\right\rangle
        =
        \langle \beta,\alpha\rangle
        \qquad
        \text{for all }\alpha,\beta\in\Phi.
    \]
    
    Therefore \(\varphi\) is an isomorphism of root systems. Hence the Cartan
    matrix determines the root system uniquely up to isomorphism.
    \end{proof}
